\documentclass[11pt]{article}
\usepackage[margin=1in]{geometry}
\usepackage{amsmath,amssymb,amsthm}
\usepackage[T1]{fontenc}
\usepackage{lmodern}
\usepackage{xcolor}
\usepackage{listings}
\usepackage{booktabs}
\usepackage{longtable}
\usepackage{array}
\usepackage[hidelinks]{hyperref}
\usepackage{url}

\newtheorem{theorem}{Theorem}
\newtheorem{lemma}[theorem]{Lemma}
\newtheorem{corollary}[theorem]{Corollary}
\newtheorem{proposition}[theorem]{Proposition}
\theoremstyle{definition}
\newtheorem{definition}[theorem]{Definition}
\theoremstyle{remark}
\newtheorem{remark}[theorem]{Remark}

% \lean{...} / \decl{...}: a Lean identifier. Every argument that looks like a Lean name is
% checked by namecheck.py against the two source trees and the pinned Mathlib.
\makeatletter
\def\UrlBreaks{\do\.\do\_\do\'}
\makeatother
\DeclareUrlCommand\leanu{\def\UrlFont{\ttfamily\upshape}}
\DeclareRobustCommand\lean[1]{\ifmmode\text{\leanu{#1}}\else\leanu{#1}\fi}
\DeclareRobustCommand\decl[1]{\lean{#1}}
% \code{...}: Lean text that is not a single identifier (not name-checked).
\newcommand{\code}[1]{\texttt{#1}}
% \ext{...}: a Lean name from a third-party development (breakable, not name-checked).
\DeclareUrlCommand\ext{\def\UrlFont{\ttfamily\upshape}}

\newcommand{\RR}{\mathbb{R}}

% version-2 commits: casey=8e9058891620926763871d85a204a7a980354f21 fuss=c5a5915b223e5f7fce56e5c5dcaba450ef36ad48
\newcommand{\commitcasey}{a83c43b8d7c8b1b98c213366cfa27eec8b8ad6c3}
\newcommand{\commitfuss}{671efdc610147022ef66e9794988ba72114f85a4}
\expandafter\def\csname commitof:casey\endcsname{\commitcasey}
\expandafter\def\csname commitof:fuss\endcsname{\commitfuss}
\newcommand{\commitof}[1]{\csname commitof:#1\endcsname}
\expandafter\def\csname repoof:casey\endcsname{casey}
\expandafter\def\csname repoof:fuss\endcsname{fuss}
\newcommand{\repo}[1]{https://github.com/jpbald93/\csname repoof:#1\endcsname}
\newcommand{\src}[4]{\href{\repo{#1}/blob/\commitof{#1}/#2/#3\#L#4}{\texttt{#3}:#4}}

\lstdefinelanguage{lean}{
  morekeywords={theorem,lemma,def,structure,inductive,where,by,fun,noncomputable},
  sensitive=true, morecomment=[l]{--}, morecomment=[s]{/-}{-/}}
\lstset{language=lean, basicstyle=\ttfamily\small, keywordstyle=\bfseries,
  columns=fullflexible, keepspaces=true, breaklines=true, frame=single, framerule=0.3pt,
  literate={ℝ}{{$\mathbb{R}$}}1 {ℕ}{{$\mathbb{N}$}}1 {≠}{{$\neq$}}1 {↔}{{$\leftrightarrow$}}2
    {→}{{$\to$}}1 {∀}{{$\forall$}}1 {∃}{{$\exists$}}1 {∧}{{$\wedge$}}1 {∈}{{$\in$}}1
    {≤}{{$\le$}}1 {ᵥ}{{$_v$}}1 {Γ}{{$\Gamma$}}1 {ω}{{$\omega$}}1 {π}{{$\pi$}}1
    {∠}{{$\angle$}}1 {•}{{$\bullet$}}1 {⟪}{{$\langle\!\langle$}}1 {⟫}{{$\rangle\!\rangle$}}1
    {‖}{{$\|$}}1 {₀}{{$_0$}}1 {₁}{{$_1$}}1 {₂}{{$_2$}}1 {₃}{{$_3$}}1 {₄}{{$_4$}}1
    {×}{{$\times$}}1}

\title{Casey's theorem and Fuss's theorem,\\ formalised in Lean~4 with Mathlib}
\author{Josh Bald\\ \small Independent researcher, Ontario, Canada\\ \small\texttt{jpbald93@gmail.com}}
\date{30 September 2026; version~2, 22 October 2026 (adds Section~\ref{sec:l1415})}

\emergencystretch=3em
\begin{document}
\maketitle

\begin{abstract}
We report two small Lean~4 formalisations, built on Mathlib, of classical theorems about circles.
The first is \emph{Casey's theorem} (1866), the generalisation of Ptolemy's theorem in which the
four points on a circle $\Gamma$ are replaced by four circles tangent to $\Gamma$ and distances are
replaced by lengths of common tangents. We prove it in every real inner product space, for circles
internally tangent to $\Gamma$ and, through a signed radius, for any mixture of internally and
externally tangent circles; with all radii zero it specialises to Mathlib's Ptolemy theorem, which
the proof reuses. The second is \emph{Fuss's theorem} (1797): if a quadrilateral is inscribed in a
circle of radius $R$ and its four side lines are tangent to a circle of radius $r$ whose centre is at
distance $d\ne R$ from the first centre, with $r>0$, then $(R^2-d^2)^2=2r^2(R^2+d^2)$. We prove it in a Euclidean
plane, together with the converse in the form of Poncelet's porism for quadrilaterals: under Fuss's
relation every non-stalling four-step chain of the tangent construction (with the hypotheses of
Theorem~\ref{thm:closes}) closes, and if the inner circle lies inside
the outer one, every point of the outer circle is a vertex of such a quadrilateral. Both
original core developments are short (about 200 and 870 lines), contain no \code{sorry} and use only Lean's
standard axioms. As far as we could determine, neither theorem had a formal proof before.
The quadrilaterals need not be convex: convexity, and contact of the inner circle with the side
segments, are not proved, so the textbook forms of the converse and of Poncelet closure for convex
bicentric quadrilaterals are covered only in this narrower form. Version~2 adds Mathlib-only
restatements with non-degenerate witnesses and a comparison with independently formalised textbook
statements.
\end{abstract}

\section{Introduction}

This paper describes Lean~4 proofs of two classical theorems about configurations of circles.
Casey's theorem is listed in Mathlib's ``1000+ theorems'' file, and Fuss's theorem is the case $n=4$ of
Poncelet's closure theorem, which is also listed; neither entry has a declaration recorded
there~\cite{Mathlib1000}.

\paragraph{Casey's theorem.}
Let $\Gamma$ be a circle and let $C_1,C_2,C_3,C_4$ be circles internally tangent to $\Gamma$ at
points $T_1,T_2,T_3,T_4$ which lie on $\Gamma$ in this cyclic order. For $i\ne j$ let $t_{ij}$ be the
length of a common external tangent of $C_i$ and $C_j$; if $C_i$ and $C_j$ have centres at distance
$\delta_{ij}$ and radii $r_i,r_j$, then $t_{ij}=\sqrt{\delta_{ij}^2-(r_i-r_j)^2}$. Then
\begin{equation}\label{eq:casey}
t_{12}\,t_{34}+t_{23}\,t_{41}=t_{13}\,t_{24}.
\end{equation}
When all four circles shrink to points, $t_{ij}=|T_iT_j|$ and \eqref{eq:casey} is Ptolemy's theorem
for the cyclic quadrilateral $T_1T_2T_3T_4$. The theorem was published by John
Casey in 1866~\cite{Casey1866} and appears in his textbook~\cite{CaseySequel}.
There, and in most later accounts, it is stated together with a converse (four circles whose
tangent lengths satisfy a relation of the form \eqref{eq:casey} are tangent to a common circle or
line); we do not formalise the converse (Section~\ref{sec:l1415} shows that a direct formalisation of it without non-degeneracy hypotheses is false). A standard variant allows some of the $C_i$ to be externally
tangent to $\Gamma$, with $t_{ij}$ then the length of an \emph{internal} common tangent when $C_i$
and $C_j$ touch $\Gamma$ from different sides; our signed-radius form covers this.

\paragraph{Fuss's theorem.}
A quadrilateral is \emph{bicentric} if it has a circumcircle $\Gamma$ (centre $O$, radius $R$) and
an incircle $\omega$ (centre $I$, radius $r$). With $d=|OI|$, Nicolaus Fuss showed in
a paper presented in 1794 and printed in 1797~\cite{Fuss1797} that
\begin{equation}\label{eq:fuss}
(R^2-d^2)^2=2r^2(R^2+d^2),\qquad\text{equivalently}\qquad
\frac1{(R-d)^2}+\frac1{(R+d)^2}=\frac1{r^2}.
\end{equation}
The triangle analogue is the Euler--Chapple formula $d^2=R(R-2r)$. Poncelet's closure theorem
(1822)~\cite{Poncelet1822,Flatto} explains both: given two circles, one inside the other, start at a
point of the outer circle, draw a tangent to the inner circle, take its second intersection with
the outer circle, draw the other tangent from there, and repeat. Whether the resulting chain closes
after $n$ steps does not depend on the starting point. For $n=3$ the closure condition is the
Euler--Chapple relation and for $n=4$ it is Fuss's relation~\eqref{eq:fuss}. The triangle case,
including the Euler--Chapple relation as a closure criterion, is formalised in a companion
paper~\cite{BaldTriangles}; here we treat $n=4$.

\paragraph{Contributions.}
\begin{itemize}
\item A proof of Casey's theorem (Theorem~\ref{thm:casey}) in an arbitrary real inner product
space, for four circles internally tangent to a circle~$\Gamma$, and a signed-radius version
(Theorem~\ref{thm:casey-le}) that allows circles externally tangent to $\Gamma$. The only
nontrivial input is Mathlib's Ptolemy theorem; the new work is the tangent-length formula
$t_{ij}=|T_iT_j|\sqrt{R-r_i}\sqrt{R-r_j}/R$ (Lemma~\ref{lem:tl}).
\item A proof of Fuss's theorem (Theorem~\ref{thm:fuss}) in a Euclidean plane for quadrilaterals
that need not be convex, with the single non-degeneracy hypothesis $d\ne R$; closure of the tangent
construction after four steps under~\eqref{eq:fuss} (Theorem~\ref{thm:closes}); existence of a
bicentric quadrilateral through every point of $\Gamma$ when $\omega$ lies inside $\Gamma$
(Theorem~\ref{thm:exists}); and the resulting ``all or nothing'' statements
(Corollary~\ref{cor:aon}).
\item Lemmas that tie the definitions to their informal reading (Section~\ref{sec:correct}),
including a formal non-vacuity theorem for Fuss's theorem and a
counterexample showing that the hypothesis $d\ne R$ cannot be dropped.
\item (Version~2.) Restatements of the main theorems with Mathlib-only types, each with an explicit
non-degenerate witness, and a comparison with textbook statements formalised independently by an AI
model (Section~\ref{sec:l1415}). The forward textbook Fuss statement is covered with the additional
hypothesis $d<R$ and without the square-characterisation clause; its converse and the Poncelet statements
are covered only with side-line tangency, without proofs of convexity or contact on the side segments.
\end{itemize}

\paragraph{Availability.}
{\sloppy The two developments are separate Lake projects, published at
\url{\repo{casey}} and \url{\repo{fuss}}.
Line references in this paper are to commits \texttt{5f6252c} (Casey) and
\texttt{671efdc} (Fuss); the Casey references are also valid at the pinned commit \texttt{a83c43b},
which adds \path{Casey/Bridge.lean} and leaves \path{Basic.lean} and \path{Main.lean} unchanged.
Version~2 of this paper adds, in a new commit of each repository (Casey \texttt{8e9058891620926763871d85a204a7a980354f21}, Fuss
\texttt{c5a5915b223e5f7fce56e5c5dcaba450ef36ad48}), the Layer-14 and Layer-15 files \path{Casey/L14.lean} (174 lines),
\path{Casey/L15Classical.lean} (394), \path{Fuss/L14.lean} (354) and \path{Fuss/L15Classical.lean}
(418) of Section~\ref{sec:l1415}; that commit changes only the root import files, \path{gate.sh} and the
README, and all earlier statements and library files are unchanged. Both use Lean \texttt{v4.33.1} and Mathlib revision
\texttt{0df444a3} (tag \texttt{v4.33.1}), pinned in \texttt{lake-manifest.json}; this is the same
Mathlib as in the companion paper~\cite{BaldTriangles}. The original Casey core consists of
\path{Casey/Basic.lean} and \path{Casey/Main.lean} (200 lines, plus its original 2-line root file). The original Fuss core consists of six
files under \path{Fuss/} (866 lines): \path{Plane.lean}, \path{Forward.lean}, \path{Porism.lean},
\path{Steps.lean}, \path{Geometry.lean} and \path{Example.lean}. Two lemmas of
\path{Fuss/Geometry.lean} are copied from the triangle development \path{poncelet-triangles}
of~\cite{BaldTriangles}, with the namespace changed and a provenance note added to the
docstring (Section~\ref{sec:fuss-setting}). Supporting material (the Layer-14 and Layer-15 reports, the blind writer's statements and report, the transcript audit, and the reviews of this version with the fixes made in response) is in the folder \texttt{Casey-Fuss} of the repository \url{https://github.com/jpbald93/LEAN-verification}.\par}

{\small\begin{longtable}{@{}>{\raggedright\arraybackslash\hangindent=1.5em\hangafter=1}p{0.50\textwidth}p{0.18\textwidth}>{\raggedright\arraybackslash}p{0.26\textwidth}@{}}
\caption{Main formal statements. Links point to the declaration line at the commit listed under
Availability.}\label{tab:thms}\\
\toprule
Lean name & In this paper & Source \\ \midrule
\endfirsthead
\multicolumn{3}{@{}l}{\small\emph{Table~\ref{tab:thms}, continued}}\\
\toprule
Lean name & In this paper & Source \\ \midrule
\endhead
\bottomrule
\endfoot
\multicolumn{3}{@{}l}{\emph{Repo \texttt{casey}}}\\
\lean{Casey.tangentCentre} & Def.~\ref{def:tc} & \src{casey}{Casey}{Basic.lean}{28} \\
\lean{Casey.tangentLength} & Def.~\ref{def:tc} & \src{casey}{Casey}{Basic.lean}{33} \\
\lean{Casey.tangentCentre_zero} & Section~\ref{sec:casey-correct} & \src{casey}{Casey}{Basic.lean}{43} \\
\lean{Casey.tangentLength_zero} & Section~\ref{sec:casey-correct} & \src{casey}{Casey}{Basic.lean}{47} \\
\lean{Casey.dist_tangentCentre_center} & Lemma~\ref{lem:spec} & \src{casey}{Casey}{Basic.lean}{51} \\
\lean{Casey.dist_tangentCentre_sq_sub_sq} & Lemma~\ref{lem:key} & \src{casey}{Casey}{Basic.lean}{62} \\
\lean{Casey.tangentLength_tangentCentre} & Lemma~\ref{lem:tl} & \src{casey}{Casey}{Basic.lean}{77} \\
\lean{Casey.cospherical_of_dist_eq} & Section~\ref{sec:casey-proof} & \src{casey}{Casey}{Main.lean}{25} \\
\lean{Casey.casey_of_le} & Theorem~\ref{thm:casey-le} & \src{casey}{Casey}{Main.lean}{34} \\
\lean{Casey.casey} & Theorem~\ref{thm:casey} & \src{casey}{Casey}{Main.lean}{61} \\
\lean{Casey.tangentCentre_spec} & Lemma~\ref{lem:spec} & \src{casey}{Casey}{Main.lean}{76} \\
\lean{Casey.tangentCentre_spec_external} & Lemma~\ref{lem:spec} & \src{casey}{Casey}{Main.lean}{90} \\
\lean{Casey.casey_zero_radii_eq_ptolemy} & Corollary~\ref{cor:ptolemy} & \src{casey}{Casey}{Main.lean}{102} \\
\multicolumn{3}{@{}l}{\emph{Repo \texttt{fuss}}}\\
\lean{Fuss.parallel_of_inner_eq_zero} & Lemma~\ref{lem:plane} & \src{fuss}{Fuss}{Plane.lean}{22} \\
\lean{Fuss.eq_or_eq_or_eq_of_inner} & Lemma~\ref{lem:plane} & \src{fuss}{Fuss}{Plane.lean}{51} \\
\lean{Fuss.vertex_smul} & Lemma~\ref{lem:plane} & \src{fuss}{Fuss}{Plane.lean}{83} \\
\lean{Fuss.vertex_rel} & Lemma~\ref{lem:vrel} & \src{fuss}{Fuss}{Forward.lean}{21} \\
\lean{Fuss.not_three_on_line} & Section~\ref{sec:fuss-forward} & \src{fuss}{Fuss}{Forward.lean}{57} \\
\lean{Fuss.fuss_vec} & Section~\ref{sec:fuss-forward} & \src{fuss}{Fuss}{Forward.lean}{71} \\
\lean{Fuss.vertex_unique} & Section~\ref{sec:fuss-closure} & \src{fuss}{Fuss}{Porism.lean}{22} \\
\lean{Fuss.chain_distinct} & Section~\ref{sec:fuss-exists} & \src{fuss}{Fuss}{Porism.lean}{31} \\
\lean{Fuss.closes_vec} & Section~\ref{sec:fuss-closure} & \src{fuss}{Fuss}{Porism.lean}{59} \\
\lean{Fuss.first_tangent} & Section~\ref{sec:fuss-exists} & \src{fuss}{Fuss}{Steps.lean}{42} \\
\lean{Fuss.other_tangent} & Section~\ref{sec:fuss-exists} & \src{fuss}{Fuss}{Steps.lean}{64} \\
\lean{Fuss.second_point} & Section~\ref{sec:fuss-exists} & \src{fuss}{Fuss}{Steps.lean}{87} \\
\lean{Fuss.step} & Section~\ref{sec:fuss-exists} & \src{fuss}{Fuss}{Steps.lean}{119} \\
\lean{Fuss.isTangent_of_mem_orthRadius} & Section~\ref{sec:fuss-setting} & \src{fuss}{Fuss}{Geometry.lean}{33} \\
\lean{Fuss.IsBicentric} & Def.~\ref{def:bic} & \src{fuss}{Fuss}{Geometry.lean}{96} \\
\lean{Fuss.fuss_of_isBicentric} & Theorem~\ref{thm:fuss} & \src{fuss}{Fuss}{Geometry.lean}{105} \\
\lean{Fuss.poncelet_chain_closes} & Theorem~\ref{thm:closes} & \src{fuss}{Fuss}{Geometry.lean}{142} \\
\lean{Fuss.exists_quad_vec} & Section~\ref{sec:fuss-exists} & \src{fuss}{Fuss}{Geometry.lean}{175} \\
\lean{Fuss.exists_bicentric_of_fuss} & Theorem~\ref{thm:exists} & \src{fuss}{Fuss}{Geometry.lean}{202} \\
\lean{Fuss.fuss_iff_exists} & Corollary~\ref{cor:aon} & \src{fuss}{Fuss}{Geometry.lean}{241} \\
\lean{Fuss.fuss_iff_forall} & Corollary~\ref{cor:aon} & \src{fuss}{Fuss}{Geometry.lean}{257} \\
\lean{Fuss.exists_iff_forall} & Corollary~\ref{cor:aon} & \src{fuss}{Fuss}{Geometry.lean}{273} \\
\lean{Fuss.fuss_iff_reciprocal} & Proposition~\ref{prop:recip} & \src{fuss}{Fuss}{Geometry.lean}{280} \\
\lean{Fuss.nonvacuous} & Proposition~\ref{prop:nonvac} & \src{fuss}{Fuss}{Example.lean}{18} \\
\end{longtable}}

\section{Setting and representation}\label{sec:setting}

Throughout, $V$ is a real inner product space and $P$ a metric affine space over it:
\begin{center}\small
\code{[NormedAddCommGroup V] [InnerProductSpace $\mathbb{R}$ V] [MetricSpace P] [NormedAddTorsor V P]}
\end{center}
For Casey's theorem no dimension assumption is made. For Fuss's theorem we add
\code{[Fact (Module.finrank $\mathbb{R}$ V = 2)]}, so that $P$ is a Euclidean plane; the model plane
\lean{EuclideanSpace}~\code{$\mathbb{R}$ (Fin 2)} is used for the non-vacuity example. We write
$\angle\,a\,p\,c$ for Mathlib's unoriented angle \lean{EuclideanGeometry.angle}.

\subsection{Casey: circles as centre and signed radius}\label{sec:casey-setting}

In the Casey development a circle is not a Mathlib \lean{EuclideanGeometry.Sphere}; it is given by
its centre and a real number, and tangency is built into the choice of centre.

\begin{definition}\label{def:tc}
Let $\Gamma$ have centre $O$ and radius $R$, and let $T$ be a point.
\begin{itemize}
\item $\lean{Casey.tangentCentre}\ O\ T\ R\ r = O+(1-r/R)(T-O)$ is the centre of the circle of radius
$r$ internally tangent to $\Gamma$ at $T$.
\item $\lean{Casey.tangentLength}\ c_1\ c_2\ r_1\ r_2=\sqrt{|c_1c_2|^2-(r_1-r_2)^2}$, with
Mathlib's \lean{Real.sqrt} (which returns $0$ on negative arguments).
\end{itemize}
\end{definition}

For $0\le r<R$ and $T\in\Gamma$ the point $c=\lean{tangentCentre}\ O\ T\ R\ r$ lies on the segment
$OT$ at distance $R-r$ from $O$ and at distance $r$ from $T$, so the circle with centre $c$ and
radius $r$ touches $\Gamma$ from inside at $T$ (Lemma~\ref{lem:spec}). For $r<0$ the same formula
puts $c$ on the ray $OT$ beyond $T$, at distance $R+|r|$ from $O$ and $|r|$ from $T$: the circle of
radius $|r|$ with centre $c$ touches $\Gamma$ from outside at $T$. So one real parameter $r\le R$
encodes both kinds of tangency, and we call it a \emph{signed radius}. With signed radii
$r_i,r_j$, the quantity $|c_ic_j|^2-(r_i-r_j)^2$ is the squared length of the common external
tangent when $r_i,r_j$ have the same sign, and of the common internal tangent (for which the
relevant radius sum is $|r_i|+|r_j|=r_i-r_j$) when $r_i\ge0\ge r_j$. This is the classical
convention for the mixed form of Casey's theorem.

\paragraph{Cyclic order.}
Following Mathlib's Ptolemy theorem, the cyclic order of $T_1,\dots,T_4$ on $\Gamma$ is expressed by
a point $p$ with $\angle\,T_1\,p\,T_3=\pi$ and $\angle\,T_2\,p\,T_4=\pi$: the diagonals $T_1T_3$ and
$T_2T_4$ meet at a point strictly inside both. For four \emph{distinct} points on a circle this is the usual meaning of ``$T_1,T_2,T_3,T_4$ in this
cyclic order'' (up to reversal); each angle hypothesis is strict betweenness
(\lean{EuclideanGeometry.angle_eq_pi_iff_sbtw}). The hypotheses also admit repeated adjacent
points, e.g.\ $T_1=T_2$ and $T_3=T_4$ at opposite ends of a diameter; the identities remain true there. Mathlib's theorem is
\begin{lstlisting}
theorem mul_dist_add_mul_dist_eq_mul_dist_of_cospherical {a b c d p : P}
    (h : Cospherical ({a, b, c, d} : Set P)) (hapc : ∠ a p c = π) (hbpd : ∠ b p d = π) :
    dist a b * dist c d + dist b c * dist d a = dist a c * dist b d
\end{lstlisting}
(\lean{EuclideanGeometry.mul_dist_add_mul_dist_eq_mul_dist_of_cospherical}, by Candales and Davidson),
and our statements take the same two angle hypotheses.

\subsection{Fuss: Mathlib spheres and tangent lines}\label{sec:fuss-setting}

In the Fuss development the circles are Mathlib spheres \code{$\Gamma$ $\omega$ : Sphere P}, with
$R=\Gamma.\mathrm{radius}$, $r=\omega.\mathrm{radius}$ and $d=\code{dist $\Gamma$.center $\omega$.center}$.
The formal statements are written with these expressions; we abbreviate them as $R,r,d$.
A line is \code{line[$\mathbb{R}$, A, B]}, the affine span of $\{A,B\}$. Tangency is Mathlib's
\lean{EuclideanGeometry.Sphere.IsTangent}: an affine subspace is tangent to $\omega$ if it contains
a point $T\in\omega$ and lies in the tangent space $\omega.\lean{orthRadius}\ T$, the affine
subspace through $T$ orthogonal to the radius at~$T$.

\begin{definition}\label{def:bic}
$\lean{Fuss.IsBicentric}\ \Gamma\ \omega\ A\ B\ C\ D$ holds if $A,B,C,D\in\Gamma$ are pairwise
distinct and the four side lines $AB$, $BC$, $CD$, $DA$ are tangent to $\omega$:
\begin{lstlisting}
def IsBicentric (Γ ω : Sphere P) (A B C D : P) : Prop :=
  A ∈ Γ ∧ B ∈ Γ ∧ C ∈ Γ ∧ D ∈ Γ ∧
    A ≠ B ∧ B ≠ C ∧ C ≠ D ∧ D ≠ A ∧ A ≠ C ∧ B ≠ D ∧
    ω.IsTangent line[ℝ, A, B] ∧ ω.IsTangent line[ℝ, B, C] ∧ ω.IsTangent line[ℝ, C, D] ∧
    ω.IsTangent line[ℝ, D, A]
\end{lstlisting}
\end{definition}

This definition does not ask for convexity, nor that the points of tangency lie on the side
segments. It therefore includes crossed (self-intersecting) quadrilaterals whose four side lines
touch~$\omega$. Section~\ref{sec:correct} discusses this choice.

The proofs work in vectors relative to the centre $I$ of $\omega$. A point $X$ becomes
$x=X-I$; the centre of $\Gamma$ becomes $o=O-I$, so $\|o\|=d$; membership $X\in\Gamma$ becomes
$\|x-o\|^2=R^2$, and for $T\in\omega$ the condition $X\in\omega.\lean{orthRadius}\ T$ becomes
$\langle x,t\rangle=r^2$ (\lean{Fuss.mem_orthRadius_iff}). The translation lemmas
\lean{Fuss.vec_of_mem}, \lean{Fuss.mem_of_vec} and \lean{Fuss.vec_of_mem_center} are one-line
rewrites. In the other direction, \lean{Fuss.isTangent_of_mem_orthRadius} says that in a plane, for
$T\ne$ the centre of $\omega$, the line through two distinct points of the tangent space at $T$ is
tangent to $\omega$; it and
\lean{Fuss.exists_mem_sphere} (a circle of positive radius in a nontrivial space has a point) are
the two lemmas copied from~\cite{BaldTriangles}.

\section{Casey's theorem}\label{sec:casey}

\subsection{Statements}

\begin{theorem}[Casey]\label{thm:casey}
Let $R>0$, let $T_1,\dots,T_4$ be points with $|T_iO|=R$, let $p$ be a point with
$\angle\,T_1\,p\,T_3=\pi$ and $\angle\,T_2\,p\,T_4=\pi$, and let $0\le r_i<R$. Put
$c_i=\lean{tangentCentre}\ O\ T_i\ R\ r_i$ and $t_{ij}=\lean{tangentLength}\ c_i\ c_j\ r_i\ r_j$. Then
$t_{12}t_{34}+t_{23}t_{41}=t_{13}t_{24}$.
\end{theorem}

\begin{lstlisting}
theorem casey {O T₁ T₂ T₃ T₄ p : P} {R r₁ r₂ r₃ r₄ : ℝ} (hR : 0 < R)
    (hr₁ : 0 ≤ r₁ ∧ r₁ < R) (hr₂ : 0 ≤ r₂ ∧ r₂ < R) (hr₃ : 0 ≤ r₃ ∧ r₃ < R)
    (hr₄ : 0 ≤ r₄ ∧ r₄ < R)
    (h₁ : dist T₁ O = R) (h₂ : dist T₂ O = R) (h₃ : dist T₃ O = R) (h₄ : dist T₄ O = R)
    (h₁₃ : ∠ T₁ p T₃ = π) (h₂₄ : ∠ T₂ p T₄ = π) :
    tangentLength (tangentCentre O T₁ R r₁) (tangentCentre O T₂ R r₂) r₁ r₂ *
        tangentLength (tangentCentre O T₃ R r₃) (tangentCentre O T₄ R r₄) r₃ r₄ +
      tangentLength (tangentCentre O T₂ R r₂) (tangentCentre O T₃ R r₃) r₂ r₃ *
        tangentLength (tangentCentre O T₄ R r₄) (tangentCentre O T₁ R r₁) r₄ r₁ =
      tangentLength (tangentCentre O T₁ R r₁) (tangentCentre O T₃ R r₃) r₁ r₃ *
        tangentLength (tangentCentre O T₂ R r₂) (tangentCentre O T₄ R r₄) r₂ r₄
\end{lstlisting}

\begin{theorem}[Casey, signed radii]\label{thm:casey-le}
The conclusion of Theorem~\ref{thm:casey} holds under the same hypotheses with $0\le r_i<R$ replaced by
$r_i\le R$ (\lean{Casey.casey_of_le}). In particular it holds for any combination of circles
internally tangent ($0\le r_i<R$) and externally tangent ($r_i<0$) to $\Gamma$.
\end{theorem}

Theorem~\ref{thm:casey} is proved from Theorem~\ref{thm:casey-le} in one line. The value $r_i=R$ is
allowed in Theorem~\ref{thm:casey-le}; then $c_i=O$ and the ``circle'' is $\Gamma$ itself, every
$t_{ij}$ involving $i$ is $0$ by Lemma~\ref{lem:tl}, and the identity reduces to $0=0$.

\begin{corollary}[Ptolemy]\label{cor:ptolemy}
Under the hypotheses on $R$, $T_i$ and $p$ of Theorem~\ref{thm:casey},
$|T_1T_2|\,|T_3T_4|+|T_2T_3|\,|T_4T_1|=|T_1T_3|\,|T_2T_4|$
(\lean{Casey.casey_zero_radii_eq_ptolemy}).
\end{corollary}

This is obtained by specialising \lean{Casey.casey} to $r_1=\dots=r_4=0$ and simplifying with
\lean{Casey.tangentCentre_zero} ($c_i=T_i$) and \lean{Casey.tangentLength_zero} ($t_{ij}=|T_iT_j|$).
It is of course weaker than the Mathlib theorem used in the proof; its purpose is to show that the
formal statement contains Ptolemy's theorem.

\subsection{Proof}\label{sec:casey-proof}

The proof follows the classical one and has three steps.

\begin{lemma}[key identity]\label{lem:key}
If $R\ne0$ and $|T_1O|=|T_2O|=R$, then for all real $r_1,r_2$, with $c_i=\lean{tangentCentre}\ O\ T_i\ R\ r_i$,
\[
|c_1c_2|^2-(r_1-r_2)^2=\frac{|T_1T_2|^2\,(R-r_1)(R-r_2)}{R^2}.
\]
\end{lemma}

\begin{proof}[Proof sketch]
Put $u=T_1-O$, $v=T_2-O$, $a=1-r_1/R$ and $b=1-r_2/R$, so $c_1-c_2=au-bv$ and $\|u\|=\|v\|=R$.
Expand $\|au-bv\|^2$ and $\|u-v\|^2$ with \lean{norm_sub_sq_real}, substitute the norms,
clear denominators (\code{field\_simp}) and close with \code{ring}. No hypothesis on the signs of
$r_i$ or on the dimension is used.
\end{proof}

\begin{lemma}[tangent length]\label{lem:tl}
If $R>0$, $r_1,r_2\le R$ and $|T_1O|=|T_2O|=R$, then
\[
t_{12}=\lean{tangentLength}\ c_1\ c_2\ r_1\ r_2=\frac{|T_1T_2|\,\sqrt{R-r_1}\,\sqrt{R-r_2}}{R}.
\]
\end{lemma}

\begin{proof}[Proof sketch]
By Lemma~\ref{lem:key} the argument of the square root is the square of the right-hand side, which
is non-negative because $R-r_i\ge0$ and $R>0$; so \lean{Real.sqrt_sq} applies.
\end{proof}

The lemma is where the hypothesis $r_i\le R$ enters: it makes $R-r_i\ge0$, so that
$\sqrt{R-r_1}\sqrt{R-r_2}=\sqrt{(R-r_1)(R-r_2)}$ and the argument of the outer square root is
non-negative. It is also the only place where the truncation of \lean{Real.sqrt} could matter,
and the lemma shows that it does not.

\begin{proof}[Proof of Theorem~\ref{thm:casey-le}]
The four points are cospherical (\lean{Casey.cospherical_of_dist_eq}, with centre $O$ and radius
$R$), so Mathlib's Ptolemy theorem gives
$|T_1T_2||T_3T_4|+|T_2T_3||T_4T_1|=|T_1T_3||T_2T_4|$. Rewrite each of the six tangent lengths
with Lemma~\ref{lem:tl}. Writing $s_i=\sqrt{R-r_i}$, the left-hand side of \eqref{eq:casey}
becomes $(|T_1T_2||T_3T_4|+|T_2T_3||T_4T_1|)\,s_1s_2s_3s_4/R^2$ (an auxiliary identity proved by
\code{field\_simp}), and the right-hand side becomes $|T_1T_3||T_2T_4|\,s_1s_2s_3s_4/R^2$.
Ptolemy's identity finishes the proof.
\end{proof}

The original Casey core is two files and 200 lines. The main work was choosing a representation in
which Lemma~\ref{lem:key} is a polynomial identity: the circles are parametrised by their points of
tangency and signed radii, not by arbitrary centres subject to a tangency hypothesis.

\section{Fuss's theorem and Poncelet's porism for quadrilaterals}\label{sec:fuss}

\subsection{Statements}

\begin{theorem}[Fuss]\label{thm:fuss}
Let $P$ be a Euclidean plane, $\Gamma,\omega$ circles with $r>0$ and $d\ne R$. If
$\lean{IsBicentric}\ \Gamma\ \omega\ A\ B\ C\ D$, then $(R^2-d^2)^2=2r^2(R^2+d^2)$.
\end{theorem}

\begin{lstlisting}
theorem fuss_of_isBicentric [Fact (Module.finrank ℝ V = 2)] (Γ ω : Sphere P) {A B C D : P}
    (hr : 0 < ω.radius) (hd : dist Γ.center ω.center ≠ Γ.radius)
    (h : IsBicentric Γ ω A B C D) :
    (Γ.radius ^ 2 - dist Γ.center ω.center ^ 2) ^ 2 =
      2 * ω.radius ^ 2 * (Γ.radius ^ 2 + dist Γ.center ω.center ^ 2)
\end{lstlisting}

\begin{theorem}[closure]\label{thm:closes}
Let $P$ be a Euclidean plane, $r>0$, $d\ne R$, and assume Fuss's relation~\eqref{eq:fuss}. Let
$A_0,\dots,A_4\in\Gamma$ and $T_0,\dots,T_4\in\omega$ be such that $A_i$ lies on the tangents to
$\omega$ at $T_i$ and $T_{i+1}$ ($i=0,\dots,3$) and $A_4$ lies on the tangent at $T_4$, with
$T_i\ne T_{i+1}$ and $A_i\ne A_{i+1}$ for $i=0,\dots,3$. Then $A_4=A_0$ and $T_4=T_0$
(\lean{Fuss.poncelet_chain_closes}).
\end{theorem}

In words: start at $A_0$ on the tangent at $T_0$; take the \emph{other} tangent from $A_0$ (touching at
$T_1\ne T_0$) and its \emph{other} intersection $A_1\ne A_0$ with $\Gamma$; repeat. After four steps
the construction is back where it started. ``$A$ lies on the tangent at $T$'' is written
$A\in\omega.\lean{orthRadius}\ T$ with $T\in\omega$. The theorem quantifies over all such chains, so
it is independent of how the choices in the construction are made.

\begin{theorem}[existence, porism]\label{thm:exists}
Let $P$ be a Euclidean plane, $r>0$ and $d+r<R$ (so $\omega$ lies strictly inside $\Gamma$), and
assume~\eqref{eq:fuss}. Then for every $A\in\Gamma$ there are $B,C,D$ with
$\lean{IsBicentric}\ \Gamma\ \omega\ A\ B\ C\ D$ (\lean{Fuss.exists_bicentric_of_fuss}).
\end{theorem}

\begin{corollary}[all or nothing]\label{cor:aon}
Let $P$ be a Euclidean plane, $r>0$ and $d+r<R$. The following are equivalent:
(i) some bicentric quadrilateral for $\Gamma,\omega$ exists;
(ii) every point of $\Gamma$ is a vertex of one;
(iii) Fuss's relation~\eqref{eq:fuss} holds.
The three equivalences (i)$\Leftrightarrow$(iii), (ii)$\Leftrightarrow$(iii) and
(i)$\Leftrightarrow$(ii) are \lean{Fuss.fuss_iff_exists}, \lean{Fuss.fuss_iff_forall} and
\lean{Fuss.exists_iff_forall}.
\end{corollary}

\begin{proposition}\label{prop:recip}
For $0\le d<R$ and $r\ne0$, $(R^2-d^2)^2=2r^2(R^2+d^2)$ if and only if
$1/(R-d)^2+1/(R+d)^2=1/r^2$ (\lean{Fuss.fuss_iff_reciprocal}).
\end{proposition}

\subsection{Three facts about the plane}

The vector proofs of the forward relation and of closure use dimension two through the following
lemma and its consequences; the existence construction (\lean{Fuss.first_tangent}) and the affine
tangency translation (\lean{Fuss.isTangent_of_mem_orthRadius}) also use dimension two directly.

\begin{lemma}\label{lem:plane}
Let $\dim V=2$.
\begin{enumerate}
\item[(a)] If $e\ne0$ and $\langle a,e\rangle=\langle b,e\rangle=0$, then $\|a\|^2b=\langle a,b\rangle a$
(\lean{Fuss.parallel_of_inner_eq_zero}).
\item[(b)] A line meets a circle in at most two points: if $\langle a-b,n\rangle=\langle a-d,n\rangle=0$ with $n\ne0$
and $b,d$ are at the same distance from $c$ as $a$, then two of $a,b,d$ coincide
(\lean{Fuss.eq_or_eq_or_eq_of_inner}).
\item[(c)] If $\|x\|^2=\|y\|^2=r^2$, $x\ne y$ and $\langle v,x\rangle=\langle v,y\rangle=r^2$ (so $v$ is the
intersection of the tangents to the circle $\|\cdot\|=r$ at $x$ and $y$), then
$(r^2+\langle x,y\rangle)\,v=r^2(x+y)$ (\lean{Fuss.vertex_smul}), and $r^2+\langle x,y\rangle\ne0$ if
$r\ne0$ (\lean{Fuss.vertex_denom_ne}).
\end{enumerate}
\end{lemma}

Part (a) is proved with Mathlib's area form on an oriented two-dimensional space
(\lean{Orientation.inner_mul_areaForm_sub}). Further consequences are that a point lies on at most
two tangents of a circle (\lean{Fuss.tangent_points}) and that a tangent line meets a circle in at
most two points (\lean{Fuss.line_points}).

\subsection{The key identity}

Place $I$ at the origin and put $o=O-I$ and $k=R^2-\|o\|^2=R^2-d^2$. For a tangency point $x$ of $\omega$
define the \emph{shifted point} $y=kx+2r^2o$, and let
\[
F=k^2-2r^2(R^2+\|o\|^2),
\]
so that Fuss's relation is $F=0$.

\begin{lemma}[shifted tangency points]\label{lem:vrel}
Let $\dim V=2$ and $r\ne0$. If $\|x\|^2=\|y\|^2=r^2$, $x\ne y$, $\langle a,x\rangle=\langle a,y\rangle=r^2$
and $\|a-o\|^2=R^2$ (the tangents at $x$ and $y$ meet at a point $a$ of $\Gamma$), then
\[
\bigl\langle kx+2r^2o,\ ky+2r^2o\bigr\rangle=-r^2F .
\]
\end{lemma}

This is \lean{Fuss.vertex_rel}.

\begin{proof}[Proof sketch]
Write $s=r^2+\langle x,y\rangle$, which is nonzero by Lemma~\ref{lem:plane}(c). From $sa=r^2(x+y)$ one
gets $s^2\|a\|^2=2r^4s$ (take the squared norm) and $s\langle a,o\rangle=r^2(\langle x,o\rangle+\langle y,o\rangle)$
(take the inner product with $o$). Expanding the left-hand side of the claim and
$\|a-o\|^2=R^2$, the product of $s$ with (left side $+\,r^2F$) is an explicit polynomial linear combination of these three
equations, certified by \code{linear\_combination}; cancelling $s\ne0$ gives
the claim.
\end{proof}

The certificate was also checked independently with SymPy (\path{code/verify_lc.py},
Section~\ref{sec:numerics}). The point of the lemma is that the right-hand side depends only on
$R,r,d$: the shifted tangency points of the two sides through any vertex have the same inner product
$\mu=-r^2F$.

\subsection{Fuss's theorem}\label{sec:fuss-forward}

The affine statement is reduced to the vector statement \lean{Fuss.fuss_vec}: the tangency points
$T_1,\dots,T_4$ are extracted from the four \lean{IsTangent} hypotheses, both endpoints of each side
lie in the tangent space at its tangency point (the field \code{le\_orthRadius} of \lean{EuclideanGeometry.Sphere.IsTangentAt}),
and $d\ne R$ becomes $k\ne0$ (using $R\ge0$, which follows from $A\in\Gamma$).

\begin{proof}[Proof of \code{fuss\_vec}]
Let the vertices be $a,b,c,e$ and the tangency points of $ab,bc,ce,ea$ be $x_1,\dots,x_4$.
\emph{Distinctness.} Three distinct points of $\Gamma$ cannot lie on one tangent line
(\lean{Fuss.not_three_on_line}, from Lemma~\ref{lem:plane}(b)). Since the vertices are pairwise
distinct, this gives $x_i\ne x_{i+1}$ (indices mod~4) and also $x_1\ne x_3$, $x_2\ne x_4$.
\emph{Equal inner products.} By Lemma~\ref{lem:vrel} at the four vertices, the shifted points
$y_i=kx_i+2r^2o$ satisfy $\langle y_4,y_1\rangle=\langle y_1,y_2\rangle=\langle y_2,y_3\rangle=\langle y_3,y_4\rangle=\mu$.
Since $k\ne0$, $x\mapsto kx+2r^2o$ is injective, so $y_1\ne y_3$ and $y_2\ne y_4$.
\emph{Contradiction if $F\ne0$.} Suppose $F\ne0$, so $\mu\ne0$. Then $y_2$ and $y_4$ are both
orthogonal to $y_1-y_3\ne0$, so by Lemma~\ref{lem:plane}(a), $\|y_4\|^2y_2=\langle y_4,y_2\rangle y_4$.
Taking the inner product with $y_3$ gives $\|y_4\|^2\mu=\langle y_4,y_2\rangle\mu$, hence
$\|y_4\|^2=\langle y_4,y_2\rangle$ and so $\|y_4\|^2(y_2-y_4)=0$. Also $y_4\ne0$ since
$\langle y_3,y_4\rangle=\mu\ne0$. Hence $y_2=y_4$ and $x_2=x_4$, a contradiction. So $F=0$.
\end{proof}

No hypothesis places $\omega$ inside $\Gamma$ or $I$ inside the quadrilateral. This differs from the
triangle case, where the analogous argument gives Euler's relation only when $\omega$ is the incircle
and a different relation, $d^2=R(R+2r)$, when it is an excircle~\cite{BaldTriangles}; for
quadrilaterals there is no second branch.

\subsection{Closure}\label{sec:fuss-closure}

\begin{proof}[Proof of Theorem~\ref{thm:closes} (via \code{closes\_vec})]
Now $F=0$, so by Lemma~\ref{lem:vrel} consecutive shifted points are orthogonal:
$\langle y_i,y_{i+1}\rangle=0$ for $i=0,\dots,3$. Each $y_i$ is nonzero: from $kx=-2r^2o$ and
$\|x\|=r$ one gets $k^2=4r^2\|o\|^2$, and together with $F=0$ this forces $r^2k=0$, which is excluded.
Also $\|y_i-2r^2o\|^2=k^2r^2$ for all $i$, so all $y_i$ lie on one circle centred at $2r^2o$.
\emph{Step 1.} $y_1$ and $y_3$ are both orthogonal to $y_2\ne0$, so they are parallel
(Lemma~\ref{lem:plane}(a)); since $y_0\perp y_1$ and $y_1\ne0$, also $y_0\perp y_3$.
\emph{Step 2.} Now $y_0,y_2,y_4$ are all orthogonal to $y_3$, so they lie on one line; they also lie
on the circle above. By Lemma~\ref{lem:plane}(b), two of them coincide. If $x_0=x_2$, then $a_0$ and $a_1$ would both lie on the two distinct tangents at $x_1$ and $x_2$,
forcing $a_0=a_1$ by \lean{Fuss.vertex_unique}; so $x_0\ne x_2$, and similarly $x_2\ne x_4$. Hence $x_4=x_0$.
\emph{Step 3.} The points $a_0,a_3,a_4$ of $\Gamma$ all lie on the tangent at $x_4=x_0$, which meets
$\Gamma$ in at most two points (\lean{Fuss.line_points}). The case $a_3=a_4$ is excluded by hypothesis;
the case $a_0=a_3$ would put $a_0$ on three distinct tangents (at $x_0$, $x_1$ and $x_3$), which leads
to a contradiction with \lean{Fuss.tangent_points} and \lean{Fuss.vertex_unique}. Hence $a_4=a_0$.
\end{proof}

\subsection{Existence}\label{sec:fuss-exists}

Theorem~\ref{thm:closes} assumes that a chain is given. To prove Theorem~\ref{thm:exists} one has to
construct one, and this is where $d+r<R$ is used.
\begin{itemize}
\item Every point of $\Gamma$ lies strictly outside $\omega$, and every point of $\omega$ strictly inside
$\Gamma$ (\lean{Fuss.outside_of_mem}, \lean{Fuss.inside_of_mem}).
\item From a point outside $\omega$ there is a tangent (\lean{Fuss.first_tangent}); the tangency point is
written down explicitly with Mathlib's right-angle rotation \lean{Orientation.rightAngleRotation}.
\item Given one tangency point $x$ from $a$, its reflection in the line $0a$ is a second, different
one (\lean{Fuss.other_tangent}).
\item The line through $a\in\Gamma$ and a point $x$ strictly inside $\Gamma$ meets $\Gamma$ again at an
explicit point $b\ne a$, which lies on the same tangent if $x$ is its tangency point
(\lean{Fuss.second_point}).
\end{itemize}
\lean{Fuss.step} combines the last two items into one step of the construction.
\lean{Fuss.exists_quad_vec} applies it four times from a first tangent at the given point $a$, closes
the chain with \lean{Fuss.closes_vec}, and obtains the two diagonal inequalities $a_0\ne a_2$,
$a_1\ne a_3$ from \lean{Fuss.chain_distinct}; these do not need Fuss's relation. Finally
\lean{Fuss.exists_bicentric_of_fuss} translates back to $P$ and proves the four tangencies with
\lean{Fuss.isTangent_of_mem_orthRadius}.

Corollary~\ref{cor:aon} follows: for (i)$\Rightarrow$(iii) and (ii)$\Rightarrow$(iii) use
Theorem~\ref{thm:fuss}, noting that $d+r<R$ and $r>0$ imply $d\ne R$; for (iii)$\Rightarrow$(ii) use
Theorem~\ref{thm:exists}, and for (ii)$\Rightarrow$(i) the fact that $\Gamma$ has a point
(\lean{Fuss.exists_mem_sphere}). Proposition~\ref{prop:recip} is a direct computation.

\section{Correctness of the formal statements}\label{sec:correct}

The Lean kernel checks the proofs, but not that the statements say what the informal theorems say.
We explain each hypothesis and give lemmas that support the intended reading.

\subsection{Casey}\label{sec:casey-correct}

\begin{lemma}\label{lem:spec}
Let $R>0$, $|TO|=R$ and $c=\lean{tangentCentre}\ O\ T\ R\ r$.
\begin{enumerate}
\item[(a)] If $0\le r<R$, then $|cO|=R-r$ and $|Tc|=r$ (\lean{Casey.tangentCentre_spec}).
\item[(b)] If $r\le0$, then $|cO|=R+(-r)$ and $|Tc|=-r$ (\lean{Casey.tangentCentre_spec_external}).
\end{enumerate}
For all $r$, $|cO|=|R-r|$ (\lean{Casey.dist_tangentCentre_center}).
\end{lemma}

Two circles with centres at distance $R-r$ and radii $R>r$ are internally tangent, and $T$ is a common
point; with distance $R+|r|$ they are externally tangent. So (a) and (b) say that \lean{tangentCentre}
places the circle correctly. Conversely, every circle of radius $r<R$ internally tangent to $\Gamma$
at $T$ has its centre on the segment $OT$ at distance $R-r$ from $O$, which is the point
\lean{tangentCentre}. This uniqueness statement is formalised in the Layer-14 file for Mathlib's
internal tangency \lean{EuclideanGeometry.Sphere.IsIntTangentAt} (\decl{Casey.l14_center_of_isIntTangentAt},
Section~\ref{sec:l1415}).

\paragraph{The hypotheses.}
\code{hR : 0 < R} makes $\Gamma$ a genuine circle; \code{h$_i$} put $T_i$ on $\Gamma$; \code{hr$_i$}
restrict to circles inside $\Gamma$, $r_i=0$ being allowed and meaning the point $T_i$; and the angle
hypotheses encode cyclic order exactly as in Mathlib's Ptolemy theorem. Without cyclic order the
statement is false: Ptolemy's identity fails for the order $(T_1,T_3,T_2,T_4)$, and so does Casey's
(this is the negative control in Section~\ref{sec:numerics}). The strict inequality $r_i<R$ in
Theorem~\ref{thm:casey} excludes the degenerate ``circle'' $\Gamma$ itself; Theorem~\ref{thm:casey-le}
shows that it is not needed for the identity.

\paragraph{Non-vacuity.}
By Corollary~\ref{cor:ptolemy} the hypotheses of Theorem~\ref{thm:casey} are those of Ptolemy's
theorem plus the radius bounds, and $r_i=0$ satisfies the bounds. Ptolemy's hypotheses are satisfied,
for example, by the vertices of a square inscribed in $\Gamma$ with $p=O$. This point-circle instance
is not a Lean theorem, but version~2 adds a non-degenerate witness with four distinct circles of
positive radius (\decl{Casey.l14w_casey}, Section~\ref{sec:l1415}).

\paragraph{Bridge lemmas.}
The module \code{Casey/Bridge.lean} connects the two definitions to the classical notions.
\decl{Casey.tangentCentre_mem_line} and \decl{Casey.collinear_tangentCentre} put $c$ on the line
$OT$, so $O$, $c$ and $T$ are collinear; \decl{Casey.tangentCentre_tangent} gives, for every $r$,
$|TO|=R$, $|Tc|=|r|$ and $|cO|=|R-r|$, the internal and external tangency distances.
\decl{Casey.tangentLength_sq} is $t^2=d^2-(r_1-r_2)^2$ when the right side is nonnegative, and
\decl{Casey.tangentLength_eq_dist_of_perp} shows that for a point and a circle the formula is the
length $|pX|$ of a tangent segment whose radius $cX$ is perpendicular to $pX$, with
$|pX|^2=|pc|^2-r^2$.

\paragraph{What is not claimed.}
The tangent length is defined by the formula, not as the length of a segment of a common tangent
line; that the two agree is the Pythagorean theorem in the right trapezoid formed by the centres and
the points of tangency, and is not formalised. The theorem is stated for four circles tangent to a
circle, not to a line.

\subsection{Fuss}

\paragraph{The hypotheses of Theorem~\ref{thm:fuss}.}
\code{hr : 0 < $\omega$.radius} excludes a point-circle. \code{hd : $d\ne R$} excludes the case where the
centre $I$ of $\omega$ lies on $\Gamma$, and it is needed: with $d=R$ the relation reads
$0=2r^2\cdot2R^2$, which is false for $r>0$, but bicentric quadrilaterals in the sense of
Definition~\ref{def:bic} still exist. The script \path{code/check_fuss.py} constructs one with $R=d=2$,
$r=1/2$, $O=(2,0)$ and $I=(0,0)$, with vertices in the order
$A\approx(0.0645,0.5039)$, $B\approx(3.8755,0.6946)$, $C\approx(0.0645,-0.5039)$,
$D\approx(3.8755,-0.6946)$. Exactly: with tangency vectors $x_1=(3/20,\sqrt{91}/20)$,
$x_2=(-1/40,\sqrt{399}/40)$, $x_3=(3/20,-\sqrt{91}/20)$, $x_4=(-1/40,-\sqrt{399}/40)$, the vertices are
$A_i=(x_i+x_{i+1})/\bigl(4(1/4+\langle x_i,x_{i+1}\rangle)\bigr)$ (indices mod 4). They are pairwise distinct, lie on $\Gamma$, and the four side lines are tangent to $\omega$.
In this configuration $\omega$ is not inside $\Gamma$ and the quadrilateral is crossed.
The theorem does \emph{not} assume that $\omega$ is inside $\Gamma$, that $I$ is inside the
quadrilateral, or that the quadrilateral is convex; so it applies to crossed quadrilaterals whose
side lines touch $\omega$, and Fuss's relation holds for these too. The distinctness conditions in
Definition~\ref{def:bic}, including the diagonal ones $A\ne C$ and $B\ne D$, are all used in the
proof (to get $x_1\ne x_3$ and $x_2\ne x_4$).

\paragraph{The hypotheses of Theorems~\ref{thm:closes} and \ref{thm:exists}.}
In Theorem~\ref{thm:closes} the conditions $T_i\ne T_{i+1}$ and $A_i\ne A_{i+1}$ express ``take the
other tangent'' and ``take the other intersection''. Without them the chain could stall, for instance
with $A_1=A_0$. The theorem does not require $\omega$ to be inside $\Gamma$, but outside that case a
chain may fail to exist (from a point inside $\omega$ there is no tangent). Theorem~\ref{thm:exists}
requires $d+r<R$ because the construction steps need it; this is also the classical setting of
Poncelet's porism.

\begin{proposition}[non-vacuity, \code{Fuss.nonvacuous}]\label{prop:nonvac}
In \lean{EuclideanSpace}~\code{$\mathbb{R}$ (Fin 2)} there exist circles $\Gamma,\omega$ with $r>0$,
$d+r<R$ and distinct centres such that Fuss's relation holds and every point of $\Gamma$ is a vertex of a
bicentric quadrilateral. The proof uses the witnesses $\Gamma$ with centre $(1,0)$ and radius $7$, and
$\omega$ with centre $(0,0)$ and radius $24/5$ (so $d+r=29/5<7$).
\end{proposition}

Here $d=1$ and $(49-1)^2=2304=2\cdot(24/5)^2\cdot50$; in reciprocal form,
$1/6^2+1/8^2=25/576=1/(24/5)^2$. So the hypotheses of Theorem~\ref{thm:exists} are satisfiable with
non-concentric circles, and \lean{IsBicentric} is not an empty notion.
(In the concentric case $d=0$, Fuss's relation is $R=\sqrt2\,r$: the square.)

\paragraph{What is not claimed.}
Definition~\ref{def:bic} asserts tangency of the side \emph{lines}. It does not say that $\omega$ is
the incircle of a convex quadrilateral $ABCD$. When $\omega$ is inside $\Gamma$ the quadrilateral
produced by the construction is classically convex, with the tangency points on the sides, but we have
not proved this, so Theorem~\ref{thm:exists} does not assert it. Accordingly, the Layer-15
comparison (Section~\ref{sec:l1415}) covers the textbook converse of Fuss's theorem and the textbook
Poncelet closure statements only with \lean{Fuss.IsBicentric} in place of a convex quadrilateral whose
contact points lie on the side segments.

\section{Numerical checks}\label{sec:numerics}

Both repositories contain Python scripts in \path{code/} that use the same definitions as the Lean
code. They are sanity checks of the statements, not part of the proofs.

\paragraph{Casey (\texttt{code/check\_casey.py}).}
Centres and tangent lengths are computed with the formulas of Definition~\ref{def:tc}, and four
points in cyclic order are sampled on random circles, with a fixed seed.
\begin{itemize}
\item 20\,000 configurations with $0\le r_i\le0.999R$ (about 20\% with one radius set to $0$): worst relative residual
of \eqref{eq:casey} $7.5\cdot10^{-12}$.
\item 20\,000 configurations with signed radii $-2R\le r_i\le0.999R$: worst relative residual $5.2\cdot10^{-11}$.
\item \emph{Negative control:} 20\,000 further configurations, relabelled in the non-cyclic order $(T_1,T_3,T_2,T_4)$,
violate \eqref{eq:casey} by more than $10^{-6}$ in all 20\,000 cases (smallest relative
residual $1.1\cdot10^{-5}$).
\item All radii $0$ (Ptolemy): residual $0$.
\end{itemize}

\paragraph{Fuss (\texttt{code/check\_fuss.py}, \texttt{code/verify\_lc.py}).}
\begin{itemize}
\item For 2\,000 random $R$, $d<0.9R$ and $r$ chosen to satisfy \eqref{eq:fuss}, the tangent
construction from a random point closes after four steps; maximum error $3.6\cdot10^{-13}$.
\item \emph{Negative control:} for $R=3$, $d=1$ and $r$ perturbed by $\pm10\%$ from the Fuss value, the smallest closure gap is approximately $3.13$.
\item A polynomial shift identity related to Lemma~\ref{lem:vrel} (with shift point $-2r^2o/k$) holds on
100 random inputs,
and the non-vacuity example gives $2304=2304$.
\item Independently of the construction, 294 configurations were found by starting from four
tangency points on $\omega$ and solving for concyclic vertices (not necessarily convex); all satisfy
\eqref{eq:fuss}, with maximum relative defect $3.5\cdot10^{-13}$.
\item The $d=R$ counterexample of Section~\ref{sec:correct}: four distinct vertices, all at distance
$2$ from $O$, while \eqref{eq:fuss} reads $0=4$.
\item \path{code/verify_lc.py} checks with SymPy the polynomial \code{linear\_combination}
certificates in \lean{Fuss.vertex_rel}, in \lean{Fuss.first_tangent} and \lean{Fuss.other_tangent},
and in \lean{Fuss.fuss_iff_reciprocal}; all seven residuals expand to $0$.
\end{itemize}
The outputs are stored in \path{code/check_fuss.out} and \path{code/verify_lc.out}. The Casey script
and \path{check_fuss.py} were re-run while this paper was prepared and reproduced the stored figures
(the current \path{check_fuss.py} additionally prints the vertices of the $d=R$ example).

\section{Verification}\label{sec:verify}

Both projects were built with prebuilt Mathlib on a separate build machine; the machine that
prepared this paper did not run Lean. According to the development notes, each was rebuilt from an
empty build directory (\code{rm -rf .lake/build; lake build}): Casey completed with 2\,454 jobs and
Fuss with 8\,713 jobs, each with exit status $0$, no errors and no warnings.

A grep of the library sources (\path{Casey/}, \path{Casey.lean}, \path{Fuss/}, \path{Fuss.lean}) for
\code{sorry}, \code{admit}, \code{native\_decide}, \code{axiom}, \code{\#eval}, \code{run\_cmd},
\code{set\_option}, \code{macro}, \code{elab}, \code{syntax}, \code{notation} and \code{import Lean}
finds nothing; we repeated it on the copies of the sources used for this paper.

{\lstset{basicstyle=\ttfamily\scriptsize}\code{\#print axioms} was run from a scratch file outside the library (\path{scratch/Ax.lean}); the
output (\path{scratch/ax.out}) is, for Casey,
\begin{lstlisting}
'Casey.dist_tangentCentre_sq_sub_sq' depends on axioms: [propext, Classical.choice, Quot.sound]
'Casey.tangentLength_tangentCentre' depends on axioms: [propext, Classical.choice, Quot.sound]
'Casey.casey_of_le' depends on axioms: [propext, Classical.choice, Quot.sound]
'Casey.casey' depends on axioms: [propext, Classical.choice, Quot.sound]
'Casey.casey_zero_radii_eq_ptolemy' depends on axioms: [propext, Classical.choice, Quot.sound]
'Casey.tangentCentre_spec' depends on axioms: [propext, Classical.choice, Quot.sound]
'Casey.tangentCentre_spec_external' depends on axioms: [propext, Classical.choice, Quot.sound]
\end{lstlisting}
and for Fuss,
\begin{lstlisting}
'Fuss.fuss_of_isBicentric' depends on axioms: [propext, Classical.choice, Quot.sound]
'Fuss.poncelet_chain_closes' depends on axioms: [propext, Classical.choice, Quot.sound]
'Fuss.exists_bicentric_of_fuss' depends on axioms: [propext, Classical.choice, Quot.sound]
'Fuss.fuss_iff_exists' depends on axioms: [propext, Classical.choice, Quot.sound]
'Fuss.fuss_iff_forall' depends on axioms: [propext, Classical.choice, Quot.sound]
'Fuss.exists_iff_forall' depends on axioms: [propext, Classical.choice, Quot.sound]
'Fuss.fuss_iff_reciprocal' depends on axioms: [propext, Classical.choice, Quot.sound]
'Fuss.nonvacuous' depends on axioms: [propext, Classical.choice, Quot.sound]
'Fuss.fuss_vec' depends on axioms: [propext, Classical.choice, Quot.sound]
'Fuss.closes_vec' depends on axioms: [propext, Classical.choice, Quot.sound]
\end{lstlisting}}
These are Lean's three standard axioms. Every theorem in Table~\ref{tab:thms} either has its own
report or is used in the proof of a theorem that has one; the definitions do not need reports.

The author's name gate \path{namecheck.py} (not included in the verification repository) checks that every Lean identifier cited
in the text exists as a declaration in the two source trees or in the pinned Mathlib, and that every
source link lands on a declaration line (for table rows, of the declaration named in that row).
Each repository also has a \path{gate.sh} that rejects forbidden tokens, builds the library,
generates its own \code{\#print axioms} queries for every declaration cited in this paper (18 for
Casey, 34 for Fuss) and fails on any axiom beyond the standard three. On the published commits both
print \code{PASS}. In the version-2 commit the gates also check the Layer-14 and Layer-15
declarations (38 for Casey, 84 for Fuss) and again print \code{PASS}. \path{tests/tamper.sh} plants six kinds of defect in scratch copies (a
\code{sorry}, a split \code{axiom}, a fake \code{\#eval} report, a macro, a weakened statement, a
wrapped extra axiom); in both repositories the gate rejects all six, each with a nonzero exit status
and the expected failure message.

The step that matters most, comparing the formal statements with the informal theorems, cannot be
automated; Section~\ref{sec:correct} is meant to help with it.
The author has checked each statement in Table~\ref{tab:thms} against its informal counterpart by hand.

\subsection{Layers 14 and 15}\label{sec:l1415}

Version~2 adds two further checks, in new files of the published repositories:
\path{Casey/L14.lean}, \path{Casey/L15Classical.lean}, \path{Fuss/L14.lean} and
\path{Fuss/L15Classical.lean}. The Layer-15 files use the namespaces \code{Casey.Classical} and
\code{Fuss.Classical}. The earlier library files are byte-identical. The gates check 38 (Casey) and 84
(Fuss) declarations, all with the standard axioms only. \path{tests/tamper.sh} rejects all six planted
defects in both repositories.

\paragraph{Layer 14: statements in Mathlib's language and non-vacuity.}
Nine main results are restated with Mathlib-only types and proved from the library theorems.
\decl{Casey.l14_casey} restates Theorem~\ref{thm:casey}. Its circles are Mathlib spheres
$C_i$ with \code{$C_i$.IsIntTangentAt $\Gamma$ $T_i$}, and the tangent lengths are written out as
$\sqrt{|c_ic_j|^2-(r_i-r_j)^2}$, because Mathlib has no tangent-length notion. No extra hypothesis is
needed, and the statement also allows $r_i=R$. \decl{Casey.l14_casey_ptolemy} restates
Corollary~\ref{cor:ptolemy}. The bridge \decl{Casey.l14_center_of_isIntTangentAt} identifies the centre
with \lean{Casey.tangentCentre}. In the Fuss file, \decl{Fuss.l14_fuss}, \decl{Fuss.l14_poncelet_chain_closes},
\decl{Fuss.l14_exists_bicentric}, \decl{Fuss.l14_fuss_iff_exists}, \decl{Fuss.l14_fuss_iff_forall},
\decl{Fuss.l14_exists_iff_forall} and \decl{Fuss.l14_fuss_iff_reciprocal} restate
Theorems~\ref{thm:fuss}--\ref{thm:exists}, Corollary~\ref{cor:aon} and Proposition~\ref{prop:recip},
with \lean{Fuss.IsBicentric} unfolded. A mechanical check of the elaborated types found no constant of
the two libraries.

Witness theorems prove every hypothesis on an explicit non-degenerate configuration:
\begin{itemize}
\item \decl{Casey.l14w_casey}: $\Gamma$ of radius $5$ and four pairwise distinct circles of radii
$1,2,3,4$ tangent at the pairwise distinct points $\pm5,\pm5i$. All six radicands are positive.
\item \decl{Fuss.l14w_fuss}, \decl{Fuss.l14w_poncelet_chain_closes} and \decl{Fuss.l14w_exists_bicentric}:
$\Gamma=((1,0),7)$, $\omega=((0,0),24/5)$ and an explicit kite with pairwise distinct vertices and
non-concentric circles.
\item \decl{Fuss.l14w_fuss_iff}: both sides of the equivalences hold for $(\Gamma,\omega)$. For a
second inner circle of radius $1$, Fuss's relation fails and no bicentric quadrilateral exists.
\item \decl{Fuss.l14w_fuss_iff_reciprocal}: $(R,d,r)=(7,1,24/5)$ satisfies both sides, and $r=1$ neither.
\end{itemize}
Proposition~\ref{prop:nonvac} is already a Mathlib-only statement and is not restated. The
externally tangent (signed-radius) reading of Theorem~\ref{thm:casey-le} is not restated, because a
Mathlib form would need mixed tangent lengths $\sqrt{\delta^2-(r_i\pm r_j)^2}$. No restatement turned
out false or needed an extra hypothesis. The kite is convex, with its contact points on the sides, but
this is not proved in Lean.

\paragraph{Layer 15: comparison with textbook statements.}
Another agent collected textbook statements of Casey's theorem (with the signed-radius form and the
converse), of Fuss's theorem and its converse, and of Poncelet's closure theorem for quadrilaterals,
with citations (Casey, Johnson, Dörrie, Yiu, Poncelet, Bos et al.). The classical sources take a
bicentric quadrilateral to be convex. An AI model was given only these statements; it saw neither the
code nor this paper, and its transcript was audited (a transcript audit in the verification repository; the
collector's transcript passed the same audit, and the collector recorded editions and pages for each source). It wrote nine Lean statements, two of them
radicand obligations, with its own definitions on \code{EuclideanSpace $\mathbb{R}$ (Fin 2)}. Among
those definitions are a convex quadrilateral whose incircle touches every side segment, tangency as a
unique common point, and cyclic order as crossing chords. A comparer then proved each statement,
copied verbatim, from the libraries, or classified it. The result: 4 covered exactly, 4 narrower,
1 false as written. No library file was changed.
\begin{itemize}
\item \emph{Covered}: Casey's theorem for disjoint internally tangent circles, the signed-radius form
with $t_{ij}=|K_iK_j|\sqrt{(R-\rho_i)(R-\rho_j)}/R$, and both radicand obligations
(\decl{Casey.Classical.cov_CaseyFuss1_a}, \decl{Casey.Classical.cov_CaseyFuss1_b} and the two
\code{\_obligation} theorems). The disjointness hypothesis is not needed.
\item \emph{Narrower}, Fuss's theorem (\decl{Fuss.Classical.cov_CaseyFuss2_a_part}): proved for the
writer's convex quadrilaterals, with the reciprocal and solved forms and $\sqrt2\,r\le R$ with
equality iff $d=0$. It needs the extra hypothesis $d<R$, which is classically automatic but not derived
(the repository's notion allows $d=R$). The clause ``$d=0$ iff the quadrilateral is a square'' is not
formalised.
\item \emph{Narrower}, converse of Fuss (\decl{Fuss.Classical.cov_CaseyFuss2_b_part}): the quadrilateral
produced is bicentric only in the sense of Definition~\ref{def:bic}. Convexity and contact points on
the side segments are not proved.
\item \emph{Narrower}, Poncelet closure (\decl{Fuss.Classical.cov_CaseyFuss3_a_part}): from one convex
bicentric quadrilateral, a quadrilateral through every point of $\Gamma$ is obtained, the
writer's tangent chain closes after exactly four steps, and there are infinitely many such
quadrilaterals. All of these quadrilaterals are bicentric only in the weaker sense: convexity and
contact on segments are not proved.
\item \emph{Narrower}, ``existence iff Fuss's relation'' (\decl{Fuss.Classical.cov_CaseyFuss3_b_part}):
the forward direction is proved as written. The converse direction and the ``through every point''
clause hold only with \lean{Fuss.IsBicentric}.
\item \emph{Writer statement false as written}: the converse of Casey's theorem
(\decl{Casey.Classical.not_classical_CaseyFuss1_c}). With $C_0$ the unit circle, $C_1$ the point $(1,0)$
and $C_2=C_3$ the point $(-1,0)$, all tangent lengths from $C_0$ vanish, so the relation holds, but no
circle or line touches all four. The textbook converse tacitly assumes distinct, non-degenerate
circles. The paper does not claim the converse.
\end{itemize}
New proofs in the Layer-15 files: for Casey, \decl{Casey.Classical.tc_eq} (a centre at the right
distances is \lean{Casey.tangentCentre}), \decl{Casey.Classical.int_dist} and
\decl{Casey.Classical.ext_dist} (the writer's internal and external tangency give the centre
distances), \decl{Casey.Classical.signed_bridge}, \decl{Casey.Classical.no_mid},
\decl{Casey.Classical.cyc_angles} (crossing chords give the angle-$\pi$ cyclic order),
\decl{Casey.Classical.casey_signed}, \decl{Casey.Classical.radicand_nonneg} and the counterexample above.
For Fuss: \decl{Fuss.Classical.tangentAt_of_side} (tangency to a side segment gives Mathlib tangency),
\decl{Fuss.Classical.ne_of_not_collinear}, \decl{Fuss.Classical.isBicentric_of_with} (the writer's convex
bicentric quadrilateral is \lean{Fuss.IsBicentric}), \decl{Fuss.Classical.fuss_of_with},
\decl{Fuss.Classical.fuss_algebra}, \decl{Fuss.Classical.fuss_of_recip},
\decl{Fuss.Classical.tangentAt_of_lineTangent}, \decl{Fuss.Classical.not_collinear_of_mem},
\decl{Fuss.Classical.not_three}, and the infinitude of~$\Gamma$.

\section{Related work}\label{sec:related}

Our searches, on 30 September 2026, covered the pinned Mathlib source and its \path{docs/100.yaml}
and \path{docs/1000.yaml}, the Lean community's list of formalised theorems, GitHub code,
repository and pull-request search, the Isabelle Archive of Formal Proofs (through its GitHub
mirror), HOL~Light, and web search. They are not exhaustive, and the claims below are ``as far as we
could determine''. A search of the Lean Zulip chat for ``Casey'', ``Fuss'' and ``bicentric'' on
30 September 2026 returned no relevant results.

\paragraph{Casey's theorem.}
Mathlib has no declaration mentioning Casey. \path{docs/1000.yaml} has an entry
\texttt{Q1370316} (``Casey's theorem'') with no declaration or author recorded~\cite{Mathlib1000},
and the theorem is not in \path{docs/100.yaml}; the Lean community's list of the 1000+ theorems shows no
formalisation. GitHub code search for ``Casey'', ``Casey theorem'' and ``casey\_theorem'' in Lean
files returned no hits; ``generalized Ptolemy'' returned only prose datasets and an unrelated notion
in SnapPy; searches of Mathlib pull requests for ``Casey'' and of repositories for ``casey theorem''
found nothing. The Isabelle AFP has an entry for Ptolemy's theorem~\cite{AFPPtolemy} but none for
Casey's, and HOL~Light has \path{100/ptolemy.ml} but nothing on Casey~\cite{HOLLight}. Our proof
reuses Mathlib's Ptolemy theorem~\cite{Mathlib} unchanged.

\paragraph{Fuss's theorem and Poncelet's porism.}
Mathlib has nothing on Fuss's theorem or bicentric polygons, and \path{docs/1000.yaml} lists
``Poncelet's closure theorem'' (\texttt{Q1785610}) with no declaration~\cite{Mathlib1000}. GitHub code
search for ``Fuss'' and ``bicentric'' in Lean files found only unrelated uses (comments, and
Fuss--Catalan numbers), and Mathlib pull-request search for ``poncelet'' and ``bicentric'' found
nothing. Web search found no formalisation of Fuss's theorem in the Isabelle AFP, HOL~Light, Coq/Rocq
or Mizar.

The closest prior work is Wang's Lean~4 repository \texttt{wwylele/Poncelet}~\cite{WangPoncelet}
(Lean \texttt{v4.34.0}, September 2026). Its theorem \ext{EuclideanGeometry.poncelet} is the general
Poncelet porism for $n$-gons and a circle inside a circle: if one proper $n$-gon is inscribed in the
outer circle and circumscribed about the inner one, then one exists through every point of the outer
circle. This is more general than our Theorem~\ref{thm:exists} in $n$, but it contains no metric
closure condition, and in particular not Fuss's relation for $n=4$ (nor our
Theorem~\ref{thm:fuss}). According to its README, its main theorem depends on
\code{native\_decide}, and hence on an axiom beyond the standard three. We read the repository but did
not build it. The companion paper~\cite{BaldTriangles} treats $n=3$, with the Euler--Chapple relation
as closure criterion, by a different method (explicit tangent constructions and Euler's formula). The
vector method used here, based on the shifted tangency points of Lemma~\ref{lem:vrel}, was not reused
from it; only two general lemmas were copied (Section~\ref{sec:fuss-setting}).

\section{Limitations and further work}\label{sec:limits}

\begin{itemize}
\item \emph{Casey: converse.} The converse (four circles whose tangent lengths satisfy a Casey
relation are tangent to a common circle or line) is substantially harder and was not attempted.
The Layer-15 writer's formalisation of it is false as written: it lacks non-degeneracy hypotheses
(\decl{Casey.Classical.not_classical_CaseyFuss1_c}).
\item \emph{Casey: degenerate and mixed cases.} Circles tangent to a line (the limit $R\to\infty$) are
not treated. The signed radius covers circles externally tangent to $\Gamma$, with $t_{ij}$ the
internal tangent between an internally and an externally tangent circle; variants that use the
internal tangent between two circles on the same side of $\Gamma$ are not treated.
\item \emph{Casey: cyclic order.} Cyclic order is expressed only through Mathlib's
diagonal-intersection hypothesis; there is no version with oriented angles or with \lean{Sbtw}.
\item \emph{Casey: representation.} In the main files circles are given by \lean{tangentCentre}, not
as Mathlib spheres with a tangency hypothesis. The Layer-14 file restates Theorem~\ref{thm:casey} with
Mathlib spheres and \lean{EuclideanGeometry.Sphere.IsIntTangentAt}, proves that every internally
tangent circle has the \lean{tangentCentre} form, and gives a non-degenerate witness; the
externally tangent (signed-radius) case of Theorem~\ref{thm:casey-le} is not restated there
(Section~\ref{sec:l1415}). The identification of \lean{tangentLength} with the length of a tangent
segment is not formalised (Section~\ref{sec:casey-correct}).
\item \emph{Fuss: convexity.} Theorem~\ref{thm:exists} produces a quadrilateral in the sense of
Definition~\ref{def:bic}; that it is convex, with $\omega$ as its incircle and the tangency points on
the sides, is classically true when $\omega$ is inside $\Gamma$ but is not proved. For the same reason
the textbook converse and Poncelet statements of Layer~15 are proved only in a narrower form, and the
forward textbook statement needs the extra hypothesis $d<R$ (Section~\ref{sec:l1415}).
\item \emph{Fuss: range of the converse.} Closure (Theorem~\ref{thm:closes}) needs only $d\ne R$, but
existence needs $d+r<R$.
\item \emph{Fuss: packaging.} The construction is stated with explicit tangency points
(\lean{orthRadius}), not with Mathlib's \lean{EuclideanGeometry.Sphere.tangentsFrom} and
\lean{EuclideanGeometry.Sphere.secondInter} as in the triangle development~\cite{BaldTriangles}.
\item \emph{Fuss: $n$-gons.} Only $n=4$ is treated. The shifted-point argument is specific to
quadrilaterals; closure conditions for larger $n$ (Fuss and Steiner gave several) would need new ideas.
\item \emph{Duplication.} Two lemmas are copied from the triangle development; a shared library or
upstreaming to Mathlib would remove this.
\end{itemize}

\section*{Acknowledgements and use of AI}
The formal proofs, the numerical scripts and a draft of this paper were prepared with the help of a
large-language-model AI assistant. The author compared each formal statement in
Table~\ref{tab:thms} with its informal counterpart by hand. Every proof is checked by the Lean
kernel, and the author takes responsibility for all content. The Layer-14 and Layer-15 files of
version~2 (Section~\ref{sec:l1415}) were also written by AI agents, the Layer-15 textbook statements
by a writer that saw neither the code nor this paper.

\begin{thebibliography}{99}
\bibitem{AFPPtolemy} L.~Bulwahn, Ptolemy's theorem, \emph{Archive of Formal Proofs}, 2016,
\url{https://www.isa-afp.org/entries/Ptolemys_Theorem.html}.
\bibitem{BaldTriangles} J.~Bald, Carnot, Viviani, De~Gua, the Japanese theorem and Poncelet's porism
for triangles, formalised in Lean~4, preprint, 2026. Lean sources of the Poncelet part:
\url{https://github.com/jpbald93/poncelet-triangles}.
\bibitem{Casey1866} J.~Casey, On the equations and properties: (1) of the system of circles touching
three circles in a plane; (2) of the system of spheres touching four spheres in space; (3) of the
system of circles touching three circles on a sphere; (4) of the system of conics inscribed to a conic,
and touching three inscribed conics in a plane, \emph{Proceedings of the Royal Irish Academy} 9
(1864--1866), 396--423, JSTOR~20488927.
\bibitem{CaseySequel} J.~Casey, \emph{A Sequel to the First Six Books of the Elements of Euclid},
Hodges, Figgis \& Co., Dublin, 1881 and later editions.
\bibitem{Flatto} L.~Flatto, \emph{Poncelet's Theorem}, American Mathematical Society, 2009.
\bibitem{Fuss1797} N.~Fuss, De quadrilateris quibus circulum tam inscribere quam circumscribere licet,
\emph{Nova Acta Academiae Scientiarum Imperialis Petropolitanae} 10 (for 1792, printed 1797),
103--125.
\bibitem{HOLLight} J.~Harrison, HOL~Light repository, \path{100/ptolemy.ml},
\url{https://github.com/jrh13/hol-light}, read 30 September 2026.
\bibitem{Mathlib} The mathlib Community, The Lean mathematical library, in \emph{Proceedings of the
9th ACM SIGPLAN International Conference on Certified Programs and Proofs (CPP 2020)}, 367--381.
Ptolemy's theorem: \path{Mathlib/Geometry/Euclidean/Sphere/Ptolemy.lean} (M.~Candales and
B.~Davidson), revision \texttt{0df444a3}.
\bibitem{Mathlib1000} The mathlib Community, \path{docs/1000.yaml}, Mathlib repository,
\url{https://github.com/leanprover-community/mathlib4}, revision \texttt{0df444a3}, accessed
30 September 2026.
\bibitem{Poncelet1822} J.-V.~Poncelet, \emph{Trait\'e des propri\'et\'es projectives des figures},
Bachelier, Paris, 1822.
\bibitem{WangPoncelet} W.~Wang, \emph{Poncelet}, \url{https://github.com/wwylele/Poncelet}, read
30 September 2026.
\end{thebibliography}

\end{document}
